\documentclass[12pt,a4paper]{article}
\usepackage[utf8]{inputenc}
\usepackage[brazil]{babel}
\usepackage{graphicx}
\usepackage{amssymb, amsfonts, amsmath}
\usepackage{float}
\usepackage{enumerate}
\usepackage[top=1.5cm, bottom=1.5cm, left=1.25cm, right=1.25cm]{geometry}

\DeclareMathOperator{\sen}{sen}

\begin{document}
\pagestyle{empty}

\begin{center}
  \begin{tabular}{ccc}
    \begin{tabular}{c}
      \includegraphics[scale=0.25]{../../biblioteca/imagem/brasao-de-armas-brasil} \\
    \end{tabular} & 
    \begin{tabular}{c}
      Ministério da Educação \\
      Universidade Federal dos Vales do Jequitinhonha e Mucuri \\
      Faculdade de Ciências Sociais, Aplicadas e Exatas - FACSAE \\
      Departamento de Ciências Exatas - DCEX \\
      Disciplina: Cálculo Diferencial e Integral I. \quad Semestre: 2021/1\\
      Prof. Me. Luiz C. M. de Aquino\\
    \end{tabular} &
    \begin{tabular}{c}
      \includegraphics[scale=0.25]{../../biblioteca/imagem/logo-ufvjm} \\
    \end{tabular}
  \end{tabular}
\end{center}

\begin{center}
 \textbf{Avaliação III}
\end{center}

\textbf{Instruções}
\begin{itemize}
 \item Todas as justificativas necessárias na solução de cada questão devem estar presentes nesta avaliação;
 \item As respostas finais de cada questão devem estar escritas de caneta;
 \item Esta avaliação tem um total de 35,0 pontos.
\end{itemize}

\begin{enumerate}
  \item Considere a função dada por $f(x) = x^3 - 12x + 7$. 
    Determine o que for solicitado abaixo.
  
    \begin{enumerate}
      \item \textbf{[3,0 pontos]} Determine os intervalos de crescimento ou de decrescimento de $f$.
      \item \textbf{[3,0 pontos]} Determine os intervalos nos quais a concavidade do gráfico de $f$ é para baixo
        ou para cima.
      \item \textbf{[4,0 pontos]} Esboce o gráfico de f.
    \end{enumerate}

  \item Considere a função dada por $f(x) = (x + 1)^2(x - 1)$.
    Determine o que for solicitado abaixo.
  
    \begin{enumerate}
      \item \textbf{[3,0 pontos]} Determine o máximo local de $f$.
      \item \textbf{[3,0 pontos]} Determine o mínimo local de $f$.
      \item \textbf{[4,0 pontos]} Determine o ponto de inflexão de $f$.
    \end{enumerate}
   
  \item \textbf{[7,0 pontos]} Ache a maior área do retângulo com a base inferior
    sobre o eixo $x$ e os vértices superiores na parábola $y = 12 - x^2$.

  \item \textbf{[8,0 pontos]} Determine o ponto da parábola $y = 2x^2$ que está mais próximo do ponto (0, 1). 
    
  \end{enumerate}
\end{document}
